\documentclass[11pt]{article}
\usepackage[T1]{fontenc}
\usepackage[margin=0.92in]{geometry}
\usepackage{amsmath,amssymb,amsthm,mathtools}
\usepackage{graphicx}
\usepackage{microtype}
\usepackage{booktabs}
\usepackage{enumitem}
\usepackage{xcolor}
\usepackage[colorlinks=true,linkcolor=blue!55!black,citecolor=blue!55!black,urlcolor=blue!55!black]{hyperref}

\newtheorem{theorem}{Theorem}
\newtheorem{lemma}{Lemma}
\newcommand{\R}{\mathbb R}
\newcommand{\C}{\mathbb C}
\newcommand{\T}{\mathbb T}
\newcommand{\cb}{c_{\mathrm b}}
\newcommand{\artanh}{\operatorname{artanh}}

\title{The Exact Bochnak Complexification Constant of the Real $\ell_1$ Plane\\
\large Full answer to the exact-value question in arXiv:2311.03531}
\author{Discovery run \texttt{fa\_banach\_001} -- \texttt{agent\_lane\_13}}
\date{August 17, 2026}

\begin{document}
\maketitle

\begin{center}
\fbox{\parbox{0.92\textwidth}{\textbf{Status: candidate full solution.}
The argument gives matching upper and lower bounds.  The lower bound uses the
classical weighted Siciak--Zakharyuta theorem after an explicit verification
that the relevant harmonic function glues subharmonically across the unit
circle.  Independent expert review is recommended.}}
\end{center}

\section{Source question}

Rodr\'iguez defines the Bochnak complexification constants
\[
 \cb(k,\ell_1^2(\R))
 =\sup\{\|\widetilde P\|:P\in\mathcal P({}^k\ell_1^2(\R)),\ \|P\|\le1\}
\]
and
\[
 \cb(\ell_1^2(\R))
 =\limsup_{k\to\infty}\cb(k,\ell_1^2(\R))^{1/k}.
\]
The source says that its exact value is unknown and records only
$2^{1/4}\le \cb(\ell_1^2(\R))\le\sqrt2$ \cite[p.~3]{RodriguezSource}.

\begin{figure}[ht]
 \centering
 \includegraphics[width=0.97\textwidth]{source_question.png}
 \caption{Definition, exact-value question, and published bounds in
 arXiv:2311.03531v2, PDF page 3.}
\end{figure}

\section{Result}

Let
\[
 G=\sum_{n=0}^{\infty}\frac{(-1)^n}{(2n+1)^2}
   =0.915965594177219\ldots
\]
be Catalan's constant.

\begin{theorem}\label{thm:main}
The exact Bochnak complexification constant of the real $\ell_1$ plane is
\[
 \boxed{\displaystyle
 \cb(\ell_1^2(\R))=\frac{e^{2G/\pi}}{\sqrt2}}
 =1.2668686397429209\ldots .
\]
In fact, for every $k\ge1$,
\[
 \cb(k,\ell_1^2(\R))
 \le \left(\frac{e^{2G/\pi}}{\sqrt2}\right)^k,
\]
and a sequence of real homogeneous polynomials asymptotically attains this
bound after taking $k$th roots.
\end{theorem}

\section{Proof intuition}

The upper bound is a one-complex-variable argument.  Restrict a real
homogeneous polynomial to the complex line $x+\zeta y$ and apply
Poisson--Jensen at $\zeta=i$.  Maximizing the resulting logarithmic integral
over the complex $\ell_1$ unit sphere is an elementary concavity problem; the
unique geometric maximum occurs when the two coordinates have equal modulus
and phase difference $\pi/2$.

For sharpness, rotate the real diamond into a square by
$u=x+y$, $v=x-y$.  Homogeneous polynomials then become one-variable
polynomials bounded by $\max(1,|t|)^k$ on the real line.  A Cayley transform
turns this into a compact weighted problem on $\T$ with weight
$W(w)=\max(|1-w|,|1+w|)$.  Its weighted extremal at the disk center is the
Poisson mean of $\log W$, namely $2G/\pi$.  The only subtle point is to show
that this Poisson solution is the global weighted extremal; an explicit
positive jump-density computation supplies that check.

\section{The sharp upper bound}

Let $P$ be a real $k$-homogeneous polynomial on $\ell_1^2(\R)$ with
$\|P\|\le1$, and take $z=x+iy\in\ell_1^2(\C)$ with $\|z\|_1\le1$.
The latter is precisely the Bochnak complexification used in the source
\cite[p.~3]{RodriguezSource}.  Apply Poisson--Jensen in the upper half-plane
to
\[
 F(\zeta)=\widetilde P(x+\zeta y).
\]
For real $t$, homogeneity and the real norm bound give
$|F(t)|\le\|x+ty\|_1^k$.  Therefore
\begin{equation}\label{eq:PJ}
 \log|\widetilde P(z)|
 \le \frac{k}{\pi}\int_{\R}
       \frac{\log\|x+ty\|_1}{1+t^2}\,dt.
\end{equation}
Zeros only improve the inequality; degenerate cases follow by continuity.

Multiplication of $z$ by a common unimodular scalar preserves its norm and
changes $\widetilde P(z)$ only by a unimodular factor.  On the unit sphere we
may consequently write
\[
 z=(r,se^{i\delta}),\qquad r,s\ge0,\qquad r+s=1.
\]
With $t=\tan\theta$ and $g(\theta)=|\cos\theta|$, the integral in
\eqref{eq:PJ}, divided by $k$, becomes
\begin{equation}\label{eq:J}
 J(r,\delta)=\log2+\frac1\pi\int_{-\pi/2}^{\pi/2}
 \log\bigl(rg(\theta)+s g(\theta-\delta)\bigr)\,d\theta.
\end{equation}
For fixed $\delta$, the integral is concave in $r$.  Reflection
$\theta\mapsto\delta-\theta$ shows that it is unchanged by
$r\leftrightarrow1-r$, hence it is maximized at $r=s=1/2$.

The evenness and $\pi$-periodicity of $g$ reduce the phase difference to
$\delta=2h$ with $0\le h\le\pi/4$.  Splitting the integral at the zeros of
the cosine terms yields
\begin{align*}
 J(1/2,2h)&=\log2+A(h),\\
 A(h)&=\frac2\pi\biggl[
  \left(\frac\pi2-h\right)\log\cos h+h\log\sin h\\
 &\hspace{22mm}+\int_0^{\pi/2-h}\log\cos t\,dt
  +\int_0^h\log\cos t\,dt\biggr].
\end{align*}
Differentiation cancels all four free logarithmic terms and gives
\[
 A'(h)=\frac2\pi\left[h\cot h-\left(\frac\pi2-h\right)\tan h\right].
\]
This is nonnegative exactly when
\begin{equation}\label{eq:chord}
 h\ge\frac\pi2\sin^2h.
\end{equation}
On $[0,\pi/4]$ the function $\sin^2h$ is convex, so it lies below the chord
joining $(0,0)$ and $(\pi/4,1/2)$; this is precisely \eqref{eq:chord}.
Thus the maximum occurs at $h=\pi/4$.  The standard Catalan integral
\[
 \int_0^{\pi/4}\log\cos t\,dt
 =-\frac\pi4\log2+\frac G2
\]
then gives
\begin{equation}\label{eq:Jmax}
 \max_{r,\delta}J(r,\delta)
 =\frac{2G}{\pi}-\frac12\log2.
\end{equation}
Combining \eqref{eq:PJ} and \eqref{eq:Jmax} proves, for every $k$,
\begin{equation}\label{eq:upper}
 \cb(k,\ell_1^2(\R))
 \le\left(\frac{e^{2G/\pi}}{\sqrt2}\right)^k.
\end{equation}

\section{A weighted polynomial lemma}

\begin{lemma}\label{lem:weighted}
For $k\ge1$, let
\[
 M_k=\sup\left\{|p(i)|:\deg p\le k,\quad
 |p(t)|\le\max(1,|t|)^k\quad(t\in\R)\right\},
\]
where complex coefficients are allowed.  Then
\[
 \lim_{k\to\infty}M_k^{1/k}=e^{2G/\pi}.
\]
The same limit holds when only real-coefficient polynomials are allowed.
\end{lemma}

\begin{proof}
Use the Cayley transform
\[
 w=\frac{t-i}{t+i},\qquad
 t=i\frac{1+w}{1-w},
\]
and set
\[
 q(w)=(1-w)^k p\left(i\frac{1+w}{1-w}\right).
\]
Then $\deg q\le k$, $q(0)=p(i)$, and the real-line constraint is equivalent
to
\begin{equation}\label{eq:circleweight}
 |q(w)|\le W(w)^k\quad(w\in\T),
 \qquad W(w)=\max(|1-w|,|1+w|).
\end{equation}

Put $\varphi(\theta)=\log W(e^{i\theta})$ and let $H$ be its Poisson
extension to the unit disk.  Since
\[
 W(e^{i\theta})=2\max\bigl(|\sin(\theta/2)|,
                               |\cos(\theta/2)|\bigr),
\]
symmetry and the Catalan integral give
\begin{equation}\label{eq:H0}
 H(0)=\frac1{2\pi}\int_0^{2\pi}\varphi(\theta)\,d\theta
 =\frac4\pi\int_0^{\pi/4}\log(2\cos u)\,du
 =\frac{2G}{\pi}.
\end{equation}

We next verify that $H$ is not merely a harmonic upper bound but the
weighted extremal.  Define on the plane
\[
 U(w)=\begin{cases}
 H(w),&|w|\le1,\\
 \log|w|+H(1/\overline w),&|w|\ge1.
 \end{cases}
\]
It is continuous, has logarithmic growth, and equals $\varphi$ on $\T$.
Its distributional Laplacian off $\T$ vanishes.  Across $\T$, the outward
radial derivative jump is
\[
 1-2\Lambda\varphi(\theta),
\]
where $\Lambda$ is the circular Dirichlet-to-Neumann operator.  On
$-\pi/2<\theta<\pi/2$ one has
$\varphi(\theta)=\log(2\cos(\theta/2))$, and a direct periodic
Hilbert-transform calculation gives
\begin{equation}\label{eq:density}
 1-2\Lambda\varphi(\theta)
 =\frac2\pi\frac{\artanh(\sin\theta)}{\sin\theta}\ge0.
\end{equation}
The quotient has limiting value $1$ at zero and the expression extends
$\pi$-periodically; its logarithmic endpoint singularities are integrable.
For reference, if $b=\tan(\theta/2)$, the principal-value integral reduces
to
\[
 \Lambda\varphi(\theta)
 =\frac12-\frac1\pi\artanh(b)(b+b^{-1})
 =\frac12-\frac1\pi\frac{\artanh(\sin\theta)}{\sin\theta}.
\]
Thus \eqref{eq:density} proves that $U$ is subharmonic on $\C$ and belongs
to the Lelong class.

Let $V_{\T,\varphi}$ be the weighted Siciak extremal.  The function $U$ is
admissible, so $V_{\T,\varphi}\ge H$ in the disk.  Conversely, every
Lelong-class subharmonic function bounded above by $\varphi$ on $\T$ is at
most the Poisson solution $H$ in the disk.  Hence
\begin{equation}\label{eq:VeqH}
 V_{\T,\varphi}=H\qquad\text{on the unit disk}.
\end{equation}

The weighted Siciak--Zakharyuta theorem for a compact set and a continuous
weight identifies this extremal with the logarithm of the asymptotic
weighted polynomial extremal \cite[Theorem~1.1, standard grading]{MSS}.
At zero, \eqref{eq:circleweight}, \eqref{eq:H0}, and \eqref{eq:VeqH}
therefore give
\[
 \lim_{k\to\infty}M_k^{1/k}=e^{H(0)}=e^{2G/\pi}.
\]

Finally write any admissible complex-coefficient polynomial as $p=a+ib$
coefficientwise, where $a,b$ have real coefficients.  On $\R$ both
$|a|$ and $|b|$ are bounded by $|p|$, and
\[
 \max(|a(i)|,|b(i)|)\ge\frac12|p(i)|.
\]
The constant factor disappears after taking $k$th roots.  This proves the
same asymptotic statement for real coefficients.
\end{proof}

\section{Sharpness and conclusion}

Use real linear coordinates
\[
 u=x+y,\qquad v=x-y.
\]
Then
\begin{equation}\label{eq:square}
 |x|+|y|=\max(|u|,|v|).
\end{equation}
Given a real polynomial $p_k$ supplied by Lemma~\ref{lem:weighted}, form the
real $k$-homogeneous polynomial
\[
 Q_k(u,v)=u^k p_k(v/u),
\]
interpreted by its homogeneous coefficient expansion.  For real $u,v$,
the defining bound for $p_k$ gives
\[
 |Q_k(u,v)|\le\max(|u|,|v|)^k.
\]
Thus $P_k(x,y)=Q_k(x+y,x-y)$ has norm at most one on the real $\ell_1$
unit ball by \eqref{eq:square}.

Evaluate its complexification at
\[
 z_0=(1/2,i/2),\qquad \|z_0\|_1=1.
\]
At this point
\[
 u_0=\frac{1+i}{2},\qquad v_0=\frac{1-i}{2},\qquad
 \frac{v_0}{u_0}=-i,\qquad |u_0|=2^{-1/2}.
\]
Real coefficients imply $|p_k(-i)|=|p_k(i)|$, so
\[
 \|\widetilde P_k\|
 \ge 2^{-k/2}|p_k(-i)|.
\]
Lemma~\ref{lem:weighted} yields
\[
 \limsup_{k\to\infty}\cb(k,\ell_1^2(\R))^{1/k}
 \ge\frac{e^{2G/\pi}}{\sqrt2}.
\]
Together with \eqref{eq:upper}, this proves Theorem~\ref{thm:main}.

\section{Verification and novelty audit}

\begin{enumerate}[leftmargin=*]
\item \textbf{Poisson--Jensen sign and normalization.}  For
$F(\zeta)=\zeta+i$, the formula at $i$ is an equality with value $\log2$;
this checks the half-plane kernel and the absence of a missing growth term.
\item \textbf{Geometric maximization.}  Concavity first eliminates the
coordinate moduli, and the derivative computation then eliminates the phase.
The maximizing point is exactly $(1/2,i/2)$, the point used for the lower
bound.
\item \textbf{Weighted-extremal gluing.}  A first proof draft used the
unbounded real-line weight directly.  The compact circle reformulation and
the positive density \eqref{eq:density} remove both the weak-admissibility
issue at infinity and the possible failure of harmonic gluing.
\item \textbf{Coefficient field.}  Passing from complex to real coefficients
costs at most a factor of two, which is harmless for the asymptotic constant.
\item \textbf{Numerical check.}  Independent semi-infinite LP reconnaissance
at $z_0$ gives kth roots $1.19627$, $1.22386$, $1.23090$, and $1.23595$
for $k=12,20,24,28$, respectively, increasing toward
$1.2668686\ldots$.  These computations are not used in the proof.
\end{enumerate}

A bounded novelty search on August 17, 2026 used the arXiv identifier, the
exact problem language, the decimal value, the Catalan-constant formula, and
the terms ``Bochnak complexification'' and ``$\ell_1^2$''.  The 2025 follow-up
paper \cite{Rodriguez2025} proves that infinite-dimensional $\ell_1$ has the
largest Bochnak constants, but its Proposition 2.8 still records only the
$2^{1/4}$ lower bound for $\ell_1^2$ and contains no Catalan-constant formula.
No prior resolution of this exact-value problem was found.  The main novelty
risk is therefore an uncatalogued equivalent computation in weighted
approximation theory.

\begin{thebibliography}{9}
\bibitem{RodriguezSource}
J. T. Rodr\'iguez,
\emph{A $2$-dimensional real Banach space with constant of analyticity less
than one}, arXiv:2311.03531v2 (2025).

\bibitem{Rodriguez2025}
J. T. Rodr\'iguez,
\emph{On the norm of the complexification of polynomials},
Studia Math. \textbf{282} (2025), 1--23,
\href{https://doi.org/10.4064/sm230615-15-4}{doi:10.4064/sm230615-15-4}.

\bibitem{MSS}
B. S. Magn\'usson, \`A. E. Sigur{\dh}ard\'ottir, and R. Sigur{\dh}sson,
\emph{Polynomials with exponents in compact convex sets and associated
weighted extremal functions: the Siciak--Zakharyuta theorem},
arXiv:2305.08260v3 (2024).
\end{thebibliography}

\end{document}
